% ---- pg-568 （书页 556） ----

这些公司在二维空间中布点后其点间距离 $d_{ik}^{(2)}$ 不可能与\hyperref[tab:12.6]{表 12.6} 中的原始(6 维)距离完全一致. 为了布点需要更大的灵活性, 而这种灵活性只有引进更多的维数才能得到.\hfill$\blacksquare$

\noindent\textbf{例 12.16}\quad (大学的多维标度变换)

\hyperref[tab:12.9]{表 12.9} 给出与 25 所美国大学有关的数据(见例 12.19). 这些数据中有: 新生入学的平均 SAT 得分, 新生在高中时排名前 $10\%$ 的人数百分比, 报考者的入学百分比, 学生教师比, 估计的年费用以及毕业率($\%$). 将度量多维标度变换算法用于经过标准化的大学数据, 便得到这些大学的二维表示, 结果示于\hyperref[fig:12.19]{图 12.19}. 注意私立大学在图的右边聚集, 而大型公立大学一般位于左侧.

\begin{mvfigure}
\mvinc{pg-568}{0.201,0.288,0.828,0.699}
\phantomsection\label{fig:12.19}\mvcap{图 12.19\quad 由度量多维标度变换得出的大学的二维表示}
\end{mvfigure}

对同一组数据用非度量多维标度变换得到的几何表示见\hyperref[fig:12.20]{图 12.20}. 对本例来说, 度量和非度量多维标度变换所给出的几何表示非常相似. 两种方法所给出的二维应力值都在 $10\%$ 左右.\hfill$\blacksquare$

经典的度量标度变换或主坐标分析, 等价于对主成分作图. 不同的软件程序选择特征向量的符号也不同, 因此初看起来两个解似乎是不同的. 然而, 如果对一个或两个坐标作镜面反射, 这两个解就完全一致(见文献[26]).

归纳起来, 多维标度变换法的关键目标是低维空间的图像. 无论何时, 只要多元数据可以表示为二维或三维空间中的图, 目测法就能对数据解释提供很大帮助.

当多元观测值有自然的数值, 且 $p$ 维空间中的欧氏距离 $d_{ik}^{(p)}$ 可以计算, 我们就可以通过使
